\documentclass[UTF8,zihao=-4]{ctexart}

\usepackage[a4paper,top=2.2cm,bottom=2.2cm,left=2.2cm,right=2.2cm]{geometry}
\usepackage{fontspec}
\usepackage{xcolor}
\usepackage{tabularx}
\usepackage{longtable}
\usepackage{booktabs}
\usepackage{array}
\usepackage{colortbl}
\usepackage{enumitem}
\usepackage{fancyhdr}
\usepackage{titlesec}
\usepackage{hyperref}
\usepackage{lastpage}

\IfFontExistsTF{DengXian}{
  \setCJKmainfont{DengXian}
  \setCJKsansfont{DengXian}
}{}
\IfFontExistsTF{Arial}{\setmainfont{Arial}}{}

\definecolor{deepblue}{HTML}{14213D}
\definecolor{sectionblue}{HTML}{1A365D}
\definecolor{tablehead}{HTML}{EAF1F8}
\definecolor{lightbox}{HTML}{F6F8FB}
\definecolor{linegray}{HTML}{CBD5E1}
\definecolor{textgray}{HTML}{334155}

\hypersetup{
  colorlinks=true,
  linkcolor=sectionblue,
  urlcolor=sectionblue,
  pdfauthor={Codex},
  pdftitle={S100P + NAS + OpenClaw Baseline 汇报}
}
\urlstyle{tt}

\pagestyle{fancy}
\fancyhf{}
\lhead{\small S100P + NAS + OpenClaw Baseline 汇报}
\rhead{\small 2026-05-28}
\cfoot{\small 第 \thepage\ 页 / 共 \pageref{LastPage} 页}
\renewcommand{\headrulewidth}{0.4pt}

\titleformat{\section}{\Large\bfseries\color{sectionblue}}{\thesection}{0.7em}{}
\titleformat{\subsection}{\large\bfseries\color{sectionblue}}{\thesubsection}{0.7em}{}
\titlespacing*{\section}{0pt}{1.0em}{0.5em}
\titlespacing*{\subsection}{0pt}{0.7em}{0.35em}

\setlist[itemize]{leftmargin=1.8em,itemsep=0.2em,topsep=0.2em}
\setlength{\parindent}{0pt}
\setlength{\parskip}{0.45em}
\setlength{\headheight}{14pt}
\setlength{\emergencystretch}{2em}
\renewcommand{\arraystretch}{1.23}

\newcolumntype{Y}{>{\raggedright\arraybackslash}X}
\newcommand{\status}[1]{\textbf{\textcolor{sectionblue}{#1}}}
\newcommand{\answerbox}[1]{
  \vspace{0.2em}
  \noindent
  \fcolorbox{linegray}{lightbox}{
    \begin{minipage}{0.96\linewidth}
      \small #1
    \end{minipage}
  }
  \vspace{0.3em}
}

\begin{document}

\begin{titlepage}
  \thispagestyle{empty}
  \vspace*{-2.0cm}
  \noindent\colorbox{deepblue}{\makebox[\dimexpr\textwidth-2\fboxsep\relax][l]{\hspace{0.4cm}\color{white}\bfseries S100P + TS-264C + OpenClaw}}
  \vspace{2.4cm}

  \begin{center}
    {\zihao{1}\bfseries\color{deepblue} S100P + NAS + OpenClaw\\Baseline 阶段性汇报\par}
    \vspace{0.8em}
    {\large 围绕 PC OpenClaw 类似效果与 AI NAS 功能复刻两个问题\par}
    \vspace{0.4em}
    {\normalsize 2026-05-28 \quad | \quad F:\textbackslash Project\textbackslash Digua\par}
  \end{center}

  \vspace{1.0em}
  \answerbox{
  当前已经从本地 workspace fallback 推进到 NAS-backed smoke baseline：S100P 常驻 OpenClaw Gateway，TS-264C 通过 NFS v4.1 持久化挂载到 \texttt{/mnt/nas/openclaw}，文档索引、每日摘要、日志诊断、浏览器截图、ROS bag session、dataset card、实验报告、稳定性采样和安全审计均可写入 NAS。
  }

  \section*{当前结论}
  \begin{itemize}
    \item S100P 可以覆盖 PC OpenClaw 的核心链路，但定位不是完整替代 PC，而是常驻入口、机器人边缘执行节点和数据落盘节点。
    \item NAS 的价值已经体现：日志、报告、数据集、截图、索引和恢复证据都从 PC 本地迁移到统一 workspace。
    \item AI NAS 常见的资料库、数据集管理、日志诊断、实验报告和安全审计能力，已经有可复现的 smoke baseline。
    \item 未完成项集中在长期稳定性、sandbox runtime、服务收敛、Home Assistant 真实配置、控制策略和真实业务数据填充。
  \end{itemize}

  \vfill
  {\small 证据入口：\path{docs/baseline_tracking.md} 与 \path{docs/baseline_report_2026-05-28_nas_backed_smoke.md}}
\end{titlepage}

\setcounter{page}{1}

\section{问题一：S100P 能否实现 PC OpenClaw 的类似效果}

\answerbox{
回答：能实现核心效果，但不是按完整桌面 PC 替代来定义。S100P 更适合作为低功耗常驻 Gateway、飞书入口、机器人侧工具执行节点和 NAS 数据落盘节点；PC 仍适合承担临时交互和重型开发。
}

\begin{longtable}{p{0.23\linewidth}p{0.17\linewidth}p{0.52\linewidth}}
\toprule
\rowcolor{tablehead}
能力 & 状态 & 当前证据 \\
\midrule
OpenClaw Gateway 常驻 & verified & Gateway active-listening，飞书 WebSocket ready，消息 received/dispatch complete 已观察到。 \\
飞书入口 & verified & 群聊/私聊均能触发 agent；\texttt{99991672} contact 权限为非阻断告警。 \\
联网搜索 & verified & Tavily 已作为 OpenClaw 搜索源验证。 \\
NAS workspace & verified & \texttt{169.254.110.209:/OpenClawWorkspace} 持久化挂载到 \texttt{/mnt/nas/openclaw}，重启后可写。 \\
浏览器自动化 & verified & Headless Chromium 打开本地页面，截图写入 NAS，PNG magic 校验通过。 \\
ROS2 状态查询 & verified & 可读取 node/topic/service 状态。 \\
ROS bag 采集 & doing & NAS-backed start/status/stop self-test 通过；长时间命名采集策略仍未定。 \\
稳定性采样 & doing & systemd timer 已切到 NAS 输出；当前 2 个 snapshot，仍需 168 小时样本。 \\
工具白名单 & doing & narrow plugin/runner 可用；broad exec 平台级阻断仍未完全验证。 \\
Sandbox & blocked & S100P 当前无 Docker/Podman/runc runtime。 \\
\bottomrule
\end{longtable}

\subsection{S100P + NAS 的直接帮助}

\begin{itemize}
  \item 常驻性：PC 不开机时，S100P 仍能承载 Gateway、飞书入口和采样任务。
  \item 统一落盘：所有日志、报告、ROS bag、dataset card、截图和索引都写 NAS。
  \item 断电恢复：Windows 托盘工具检查 PC 到 S100P、NAS、OpenClaw/飞书的链路；S100P 侧 NFS automount 已通过重启验证。
  \item 机器人侧闭环：S100P 靠近 ROS2/TROS 和机器人数据，NAS 做长期存储，OpenClaw 做触发、诊断和汇总。
\end{itemize}

\section{问题二：AI NAS / OpenClaw NAS 能力复刻}

\answerbox{
回答：高价 AI NAS 的核心价值不只是本地模型，而是把资料库、检索、日志分析、实验报告、数据集和自动化审计组织成长期可用的工作流。当前 S100P + TS-264C 已经复刻了其中一批基础能力。
}

\begin{longtable}{p{0.28\linewidth}p{0.48\linewidth}p{0.16\linewidth}}
\toprule
\rowcolor{tablehead}
AI NAS 常见能力 & S100P + NAS 当前实现 & 状态 \\
\midrule
统一 workspace & TS-264C NFS \texttt{/OpenClawWorkspace} 挂到 S100P & verified \\
文档索引/每日摘要 & Markdown 索引、hash、preview、deterministic daily summary & verified \\
图片 caption / 搜索 & 图片 metadata caption 与 JSONL index & doing \\
日志诊断 & 从 NAS logs 输出错误摘要、关键匹配和建议命令 & verified \\
机器人数据集管理 & ROS bag session 写 NAS，并生成 \texttt{DATASET\_CARD.md} & verified \\
实验/周报 & 从 NAS logs、reports、datasets 汇总 Markdown 报告 & verified \\
稳定性报告 & NAS-backed snapshot + summary，timer 自动采样 & doing \\
安全审计 & Gateway 暴露、NAS mount、secret scan、服务监听审计 & doing \\
设备只读状态 & Home Assistant read-only preflight 已有 & doing \\
低风险控制 & control policy/audit preflight 已有，未开放执行 & doing \\
\bottomrule
\end{longtable}

\section{关键 NAS-backed 证据}

\small
\begin{longtable}{p{0.30\linewidth}p{0.62\linewidth}}
\toprule
\rowcolor{tablehead}
项目 & 路径 \\
\midrule
Baseline roll-up & \path{/mnt/nas/openclaw/reports/baseline-status/baseline_status_20260528-184444.md} \\
Experiment report & \path{/mnt/nas/openclaw/reports/experiments/experiment_report_20260528-184444.md} \\
Document index & \path{/mnt/nas/openclaw/reports/document_index_20260528-182111.md} \\
Document daily summary & \path{/mnt/nas/openclaw/reports/daily-summary/document_daily_summary_20260528-184329.md} \\
Browser screenshot & \path{/mnt/nas/openclaw/reports/browser-smoke/browser_smoke_20260528-182111.png} \\
ROS bag dataset & \path{/mnt/nas/openclaw/robot_datasets/rosbag_session_20260528-182117/} \\
Dataset card & \path{/mnt/nas/openclaw/robot_datasets/rosbag_session_20260528-182117/DATASET_CARD.md} \\
Log diagnosis & \path{/mnt/nas/openclaw/logs/probes/log_diagnosis_20260528-181546.md} \\
Image caption index & \path{/mnt/nas/openclaw/reports/image-captions/image_caption_index_20260528-182530.md} \\
Security audit & \path{/mnt/nas/openclaw/logs/probes/security_audit_20260528-182530.md} \\
Service policy & \path{/mnt/nas/openclaw/logs/probes/service_policy_20260528-183619.md} \\
\bottomrule
\end{longtable}
\normalsize

\section{当前阻塞与后续计划}

\begin{longtable}{p{0.24\linewidth}p{0.34\linewidth}p{0.34\linewidth}}
\toprule
\rowcolor{tablehead}
事项 & 当前阻塞 & 下一步 \\
\midrule
A-010 稳定性 & 需要 168 小时样本 & timer 已在 NAS 采样，等待时间累计后生成验收 summary。 \\
A-006 Sandbox & 无 Docker/Podman/runc & 决定安装 runtime 或从第一版 baseline drop。 \\
B-003 图片 caption & 当前是 metadata caption & 决定是否接语义视觉模型。 \\
B-008 设备状态 & 缺 Home Assistant URL/token & 用户提供配置后仅调用 \texttt{GET /api/} 和 \texttt{GET /api/states}。 \\
B-009 控制动作 & 缺 allowlist 与二次确认策略 & 先创建 disabled policy，再审阅 request/approve/execute。 \\
B-010 服务收敛 & 需确认 NFS/RPC、x11vnc、iiod 是否可关 & 只在确认后执行 disable/firewall。 \\
\bottomrule
\end{longtable}

\section{汇报用结论}

当前可以汇报：S100P + NAS 已经跑通一个可恢复、可观察、可写 NAS 的 OpenClaw 常驻 baseline。它已经能覆盖 PC OpenClaw 的常驻入口、消息触发、工具执行、浏览器截图、ROS 数据采集、日志诊断和报告汇总的一部分核心能力；同时，它也开始具备 AI NAS 的资料库、日志、数据集、日报/周报和安全审计雏形。

下一阶段不应继续堆 smoke demo，而应把真实文档、真实图片、真实机器人数据和 7 天稳定性样本接入同一套 NAS-backed 报告链路，并在服务收敛和控制策略确认后再进入更接近产品化的验收。

\end{document}
